\documentclass[10pt,a4paper]{amsart}
\usepackage[margin=19mm]{geometry}
\usepackage{amsmath,amssymb,mathtools}
\usepackage[scr=rsfs,cal=euler]{mathalfa}
\usepackage{xfrac}
\usepackage[unicode,hidelinks,bookmarks=false]{hyperref}
\newcommand{\ie}{i.e.\ }
\newcommand{\cf}{cf.\ }
\newcommand{\resp}{respectively\ }
\newcommand{\Subsection}{\subsection}
\providecommand{\nobreakdash}{\nobreak\hbox{-}}
\setlength{\emergencystretch}{2em}
\multlinegap0pt
\usepackage[shortlabels]{enumitem}
\usepackage{amssymb}
\usepackage{mathtools}
\usepackage{tikz}
\usepackage{caption}
\usepackage{tcolorbox}
\newtheorem{cor}{Corollary}
\newtheorem{thm}{Theorem}
\title{The symplectic left companion of a Littlewood-Richardson-Sundaram tableau \\and the Kwon property}
\author{Olga Azenhas}
\date{Source: arXiv:2601.06930v2 (CC BY 4.0)}
\begin{document}
\maketitle

\noindent\emph{Source and licence.} The symplectic left companion of a Littlewood-Richardson-Sundaram tableau \\and the Kwon property. Version arXiv:2601.06930v2 (CC BY 4.0), distributed by arXiv under the Creative Commons Attribution 4.0 International licence, http://creativecommons.org/licenses/by/4.0/, as recorded in the arXiv licence metadata for this exact version and shown on the abstract page https://arxiv.org/abs/2601.06930v2. Full licence notice: \texttt{licenses/arxiv-2601.06930-CC-BY-4.0.txt}.\par\smallskip

\noindent\textit{Editorial scope.} Counted unit: the article's thm th:main (source0.tex lines 2292--2306). The article's complete original proof of it is delivered with it (source0.tex lines 2307--2508, one \texttt{proof} environment of 202 lines). No mathematics is written, repaired or re-derived: the statement and the proof are the article's own lines, in the article's own notation, and no retained line is substituted or re-flowed.  Retained uncounted as the source-local context the proof needs: lines 2059--2061; lines 2064--2073.  Nothing was omitted from any retained span, so the package carries no omission record.  No retained line contains a citation, so the wrapper carries no bibliography and no bibliography entry.  No retained line contains a figure environment, an image-inclusion command or drawing code, so no image is delivered and the package declares no asset dependency.  Editorial formatting only, in addition: an AMS \texttt{amsart} preamble that restates the article's own declarations and macros used by the retained lines, namely the environments \texttt{cor}, \texttt{thm} on separate counters, together with the standard packages those lines need.  The retained environments print in this document as Corollary 1, Theorem 1 in order of appearance, whereas the article prints the same items with the numbering of its own full document.  The complete native source of the article as distributed in the arXiv e-print for this exact version is bundled unchanged as \texttt{sources/192453/source0.tex}.
\begin{cor} \label{neglrs} Let $T\in LR(\lambda/\mu,\nu)$ with $\nu$ even, on the alphabet $[2n]$, and $n\ge \ell(\mu)$. $T$ does not satisfy Sundaram property if and only if $T(n+t,1)=2i+1$, for some  $t> i\ge 0$.

\end{cor}

\begin{cor} \label{cor:violation} Let $T\in LR(\lambda/\mu,\nu)$ with $\nu$ even, on the alphabet $[2n]$, and $n\ge \ell(\mu)$. Then
\begin{enumerate}
\item $T(n+t,1)=even< 2t$, for some $t\ge 1$  only if $T(n+s,1)=odd< 2s+1$, for some $1\le s<t$.

\item If $s\ge 1$ is such that $(n+s,1)$ is the first cell, seen from the top, where $T(n+s,1)=odd<2s+1$ then   $$T(n+1,1)\ge 2,\dots,T(n+s-2,1)\ge 2(s-2),  T(n+s-1,1)= 2(s-1),$$ 
    and $$T(n+s,1)=2(s-1)+1, T(n+s+1,1)=2s<2(s+1).$$

\item  $T(n+s,1)=odd< 2s+1$ for some $s\ge 1$, only if $T(n+t,1)=even< 2t$, for some  $1\le s<t$.
\end{enumerate}
\end{cor}

\begin{thm}(Main result)\label{th:main} Let $T\in LR(\lambda/\mu,\nu)$ with $\nu$ even, on the alphabet $[2n]$, and $ \ell(\mu)\le n$.
 $T$ does not satisfy the Sundaram property if and only if $G_\mu(T)$ is not symplectic.

Moreover,  in this case, there exists a unique $t\ge 0$ such that
\begin{enumerate}

\item  ${n+t+1}$  is  the minimal row number of $T$ where the  Sundaram property  violation occurs.

\item
  \begin{align}&T(n+t,1)=2t,\; T(n+t+1,1)= 2t+1,\; T(n+t+2,1)=2(t+1),\\
 &T(n+1,1)\ge 2,\, T(n+2,1)\ge 4,\dots, T(n+t+1-2,1)\ge 2(t+1-2).
 \end{align}
 \item  the maximal row    of $G_\mu(T)$ where a symplectic violation  occurs is among the bottom most $t+1$ cells,  $(\ell(\mu),1)$,  $(\ell(\mu)-1,1),\dots$,  $ (\ell(\mu)-t,1) $ of the first column of $G_\mu(T)$.
\end{enumerate}
\end{thm}

\begin{proof} Let $ G=G_\mu(T)$ and recall $\ell(\lambda)-n\le n$. From Corollary \ref{neglrs}, let $t+1\ge 1$ be  minimal such that $T(n+t+1,1)=2i+1$ for some  $0\le i<t+1$.

For readability we start to spelling out the cases $t+1=1,2$.

If $\mathbf{t+1=1}$, $T(n+1,1)=1$  and $n=\ell(\mu)$. Since $\nu$ is even,  $T(n+2,1)=2$. Hence $$2\le \ell(\lambda)-n\le n\Rightarrow n=\ell(\mu)\ge 2.$$ We show that $T(n+1,1)=1$ means a symplectic violation in the cell $(\ell(\mu),1)$ of $G_\mu(T)$.

One has, $\ell(\mu)\ge 2$, and $T(\ell(\mu)+1,1)=1$,   $T(\ell(\mu)+2,1)=2$ is equivalent to
\begin{align}&\ell(\mu)=\ell(\mu^{2n})=\ell(\mu^{2n-1})=\ell(\mu^{2n-2})\ge\ell(\mu^{2n-3})\\
&\Leftrightarrow \\
&G(\ell(\mu),1)\le 2\ell(\mu)-2=2(\ell(\mu)-1)\Leftrightarrow G(\ell(\mu),1)\ngeq 2\ell(\mu)-1.
\end{align}
Hence,  $G_\mu(T)$ is not symplectic with bottom most symplectic violation in the cell $(\ell(\mu),1)$.


If $\mathbf{t+1=2}$,  by definition of $t+1$, $T(n+2,1)=3$ and $T(n+1,1)=2$. On the other hand, $\nu$ is even,  so $T(n+3,1)=4$ and $3\le \ell(\lambda)-n\le n$ and thus $n\ge 3$, and
$$T(n+1,1)=2, \;T(n+2,1)=3, T(n+3,1)=4.$$

We show that $T(n+2,1)=3$ means a symplectic violation in a cell of the first column of $G_\mu(T)$.

\medskip
\emph{Case $n=\ell(\mu)\ge 3$}:   $T(\ell(\mu)+1,1)=2$, $T(\ell(\mu)+2,1)=3$,
$T(\ell(\mu)+3,1)=4$.

This translates to

\begin{align*}&\ell(\mu)=\ell(\mu^{2n})>\ell(\mu^{2n-1})=\ell(\mu^{2n-2})=\ell(\mu^{2n-3})=\ell(\mu^{2n-4})\ge\cdots \\
&\Leftrightarrow \\
&G(\ell(\mu),1)=2n\ge 2n-1,\;\;
G(\ell(\mu)-1,1)\le 2n-4=2(\ell(\mu)-1)-2 \ngeq 2(\ell(\mu)-1)-1
\end{align*}
%Hence, $ G(\ell(\mu),1)=2n\ge 2n-1$ and  $ G(\ell(\mu)-1,1)\ngeq 2(\ell(\mu)-1)-1$,
So
$G$ is not symplectic with bottom most symplectic violation in the cell $(\ell(\mu)-1,1)$.

\medskip

\emph{Case  $n=\ell(\mu)+1\ge 3$} in which case and $\ell(\mu)\ge 2$, $T(\ell(\mu)+1,1)=1$, $T((\ell(\mu)+1)+1,1)=2$,
$T((\ell(\mu)+1)+2,1)=3$, $T((\ell(\mu)+1)+3,1)=4$.

This translates to

\begin{align*}&\ell(\mu)=\ell(\mu^{2n})=\ell(\mu^{2n-1})=\ell(\mu^{2n-2}=\ell(\mu^{2n-3})=\ell(\mu^{2n-4}\ge\cdots\nonumber \\
&\Leftrightarrow\\
 &G(\ell(\mu),1)\le 2n-4=2(\ell(\mu)+1)-4=2\ell(\mu)-2\ngeq 2\ell(\mu)-1.
\end{align*}
Hence,  $G_\mu(T)$ is not symplectic with bottom most symplectic violation in the cell $(\ell(\mu),1)$.
\medskip

\iffalse
If $\mathbf{t=3}$, by definition of $t$, $T(n+3,1)=5$, $T(n+2,1)=4$ and $T(n+1,1)=2 \mbox{ or $3$ }$. On the other hand, $\nu$ is even,  so $T(n+4,1)=6$ and $4\le \ell(\lambda)-n\le n$ and thus $n\ge 4$. Therefore, either

\begin{align}&T(n+1,1)=2, T(n+2,1)=4, T(n+3,1)=5, T(n+4,1)=6  \label{one}\\
&\mbox{ or}\nonumber\\
&T(n+1,1)=3, T(n+2,1)=4, T(n+3,1)=5, T(n+4,1)=6.\label{two}
\end{align}

In \eqref{one}, we may have $n=\ell(\mu)$, or $n=\ell(\mu)+1$ with $T(n,1)=1$.

In \eqref{two}, we may have $n=\ell(\mu)$, or $n=\ell(\mu)+1$ with $T(n,1)=2$, or $n=\ell(\mu)+2$ with $T(n-1,1)=1$, $T(n,1)=2$.

\medskip
\emph{Case $n=\ell(\mu)\ge 4$}:

$(i)$ $T(n+1,1)=2, T(n+2,1)=4, T(n+3,1)=5, T(n+4,1)=6$;

\begin{align*}&\ell(\mu)=\ell(\mu^{2n})>\ell(\mu^{2n-1})=\ell(\mu^{2n-2})>\nonumber
\\
&>\ell(\mu^{2n-3})=\ell(\mu^{2n-4})=\ell(\mu^{2n-5})=\ell(\mu^{2n-6})\ge\cdots\nonumber \\
&\Leftrightarrow\\
&G(\ell(\mu),1)=2n\ge 2n-1,\; G(\ell(\mu)-1,1)=2n-2\ge 2(n-1)-1,\\
&\;G(\ell(\mu)-2,1)\le 2n-6=2\ell(\mu)-6=2(\ell(\mu)-2)-2\ngeq 2(\ell(\mu)-2)-1.
\end{align*}

%Hence, $ G(\ell(\mu),1)=2n\ge 2n-1$, $ G(\ell(\mu)-1,1)=2n-2\ge 2(n-1)-1$ and $G(\ell(\mu)-2,1)\ngeq 2(\ell(\mu)-2)-1$ and
$G$ is not symplectic.

\medskip
$(ii)$ $T(n+1,1)=3, T(n+2,1)=4, T(n+3,1)=5, T(n+4,1)=6$

\begin{align*}&\ell(\mu)=\ell(\mu^{2n})>\ell(\mu^{2n-1})>\ell(\mu^{2n-2})=
\ell(\mu^{2n-3})=\ell(\mu^{2n-4})=\ell(\mu^{2n-5})=\ell(\mu^{2n-6}\ge\cdots\nonumber \\
&\Leftrightarrow \\
&G(\ell(\mu),1)=2n, \; G(\ell(\mu)-1,1)=2n-1,\\
&G(\ell(\mu)-2,1)\le 2n-6=2\ell(\mu)-6\ngeq 2(\ell(\mu)-2)-1.
\end{align*}
%Hence, $ G(\ell(\mu),1)=2n$, $ G(\ell(\mu)-1,1)=2n-1$, and $ G(\ell(\mu)-2,1)\ngeq 2(\ell(\mu)-2)-1$ and
$G$ is not symplectic.

\medskip
\emph{Case $n=\ell(\mu)+1\ge 4\Rightarrow\ell(\mu)\ge 3 $}

$(i)$ $T(\ell(\mu)+1,1)=1$, $T(\ell(\mu)+2,1)=2$,
$T(\ell(\mu)+3,1)=4$, $T(\ell(\mu)+4,1)=5$, $T(\ell(\mu)+5,1)=6$

 \begin{align*}&\ell(\mu)=\ell(\mu^{2n})=\ell(\mu^{2n-1})=\ell(\mu^{2n-2}>
\ell(\mu^{2n-3})=\ell(\mu^{2n-4})=\ell(\mu^{2n-5})=\ell(\mu^{2n-6})\ge\cdots\nonumber \\
&\Leftrightarrow\\
&G(\ell(\mu),1)=2n-2=2(\ell(\mu)+1)-2=2\ell(\mu)\ge 2\ell(\mu)-1\\
 &G(\ell(\mu)-1,1)\le 2n-6=2\ell(\mu+1)-6=2\ell(\mu)-4\ngeq 2(\ell(\mu)-1)-1.
\end{align*}
%Hence, $ G(\ell(\mu),1)=2n-2=2(\ell(\mu)+1)-2=2\ell(\mu)\ge 2\ell(\mu)-1$, and $ G(\ell(\mu)-1,1)\ngeq 2(\ell(\mu)-1)-1$ and
$G$ is not symplectic.

\medskip
$(ii)$ $T(\ell(\mu)+1,1)=2$, $T(\ell(\mu)+2,1)=3$,
$T(\ell(\mu)+3,1)=4$, $T(\ell(\mu)+4,1)=5$, $T(\ell(\mu)+5,1)=6$

 \begin{align*}&\ell(\mu)=\ell(\mu^{2n})>\ell(\mu^{2n-1})=\ell(\mu^{2n-2})=
\ell(\mu^{2n-3})=\ell(\mu^{2n-4})=\ell(\mu^{2n-5})=\ell(\mu^{2n-6})\ge\cdots\nonumber \\
&\Leftrightarrow\\
&G(\ell(\mu),1)=2n=2(\ell(\mu)+1),\\
& G(\ell(\mu)-1,1)\le 2n-6=2(\ell(\mu)+1)-6=2\ell(\mu)-4\ngeq 2(\ell(\mu)-1)-1.
\end{align*}
%Hence, $ G(\ell(\mu),1)=2n=2(\ell(\mu)+1)$,  $G(\ell(\mu)-1,1)\ngeq 2(\ell(\mu)-1)-1$ and
$G$ is not symplectic.

\medskip
\emph{Case $n=\ell(\mu)+2\ge 4\Rightarrow\ell(\mu)\ge 2$}

$T(\ell(\mu)+1,1)=1$, $T(\ell(\mu)+2,1)=2$, $T(\ell(\mu)+3,1)=3$, $T(\ell(\mu)+4,1)=4$, $T(\ell(\mu)+5,1)=5$, $T(\ell(\mu)+6,1)=6$

\begin{align*}&\ell(\mu)=\ell(\mu^{2n})=\ell(\mu^{2n-1})=\ell(\mu^{2n-2})=
\ell(\mu^{2n-3})=\ell(\mu^{2n-4})=\ell(\mu^{2n-5})=\ell(\mu^{2n-6})\ge\cdots\nonumber \\
&\Leftrightarrow\\
& G(\ell(\mu),1)\le 2n-6=2\ell(\mu+2)-6=2\ell(\mu)-2\ngeq 2\ell(\mu)-1.
\end{align*}

%Hence, $ G(\ell(\mu),1)\ngeq 2\ell(\mu)-1$ and
$G$ is not symplectic.


\fi

\bigskip

Le $\mathbf{t+1}\ge 1$ with $t\ge 0$, be  minimal such that $T(n+t+1,1)=2i+1$ for some  $0\le i\le t$. From Corollary \ref{cor:violation}, this means that
$T(n+k,1)$ does not violate the Sundaram property for $1\le k\le t$ and  $T(n+k,1)\ge 2k$, for $1\le k\le t$. Therefore, since $T(n+t,1)\ge 2t$, $T(n+t+1,1)\le 2t+1$ and
$$2\le T(n+1,1)<\cdots<T(n+t,1)<T(n+t+1,1)\le 2t+1,$$
it follows that $$T(n+t,1)=2t<T(n+t+1,1)= 2t+1$$

 On the other hand, $\nu$ is even,  so $T(n+t+2,1)=2(t+1)$. Thus  $\mathbf{n+t+1}\ge n+1$, with $t\ge 0$, is  the minimal row of $T$ where the  Sundaram property  violation occurs if and only if
 $$ T(n+1,1)<\cdots <T(n+t+1-2,1)<T(n+t,1)=2t,\; T(n+t+1,1)= 2t+1,\; T(n+t+2,1)=2(t+1),$$
 $$T(n+1,1)\ge 2\times 1,\, T(n+2,1)\ge 2\times 2,\dots, T(n+t+1-2,1)\ge 2(t+1-2).$$
 Moreover $t+2\le \ell(\lambda)-n\le n$ and thus $n\ge t+2$.

 \medskip
 %\emph{We show that if the first Sundaram condition  violation  occurs with $T(n+t+1,1)= 2t+1$ then it incurs in a symplectic violation  either %in  the cell $(\ell(\mu),1)$, $(\ell(\mu)-1,1),\dots$, or $ (\ell(\mu)-t,1) $  of the first column of $G_\mu(T)$.}



 \emph{Case $n=\ell(\mu)\ge t+1+1$}:

\medskip
Let $T(n+1,1)=2, T(n+2,1)=4,\dots, T(n+t-1),1)=2(t-1),\;$ and $T(n+t,1)=2t,\;$ $T(n+t+1,1)= 2t+1,$ $\; T(n+t+2,1)=2(t+1)$. Note $2,4,\dots, 2(t-1)$ are the first positive $t-1$ even numbers.

 This translates to
\begin{align}\label{vip}\\
&\ell(\mu)=\ell(\mu^{2n})>\nonumber\\
&\ell(\mu^{2n-1})=\ell(\mu^{2n-2})>\nonumber\\
&\ell(\mu^{2n-3})=\ell(\mu^{2n-4})>\nonumber\\
&\vdots \qquad\qquad\qquad\qquad\qquad\nonumber\\
%&\qquad\qquad\qquad\qquad\qquad\qquad\qquad >\nonumber\\
&\ell(\mu^{2n-(2(t+1-2)-1)})=\ell(\mu^{2n-2(t+1-2)})>\nonumber
\\
&\ell(\mu^{2n-(2(t+1-2)+1)})=\ell(\mu^{2n-(2(t+1-2)+2)})=\ell(\mu^{2n-(2(t+1-2)+3)})=\ell(\mu^{2n-(2(t+1-2)+4)})\ge\cdots \label{flat} \\
&\Leftrightarrow \nonumber\\
&G(\ell(\mu),1)=2n\ge 2n-1\;,  G(\ell(\mu)-1,1)=2n-2\ge 2(n-1)-1, \dots,\nonumber\\ &G(\ell(\mu)-(t+1-2),1)=2\ell(\mu)-2(t+1-2)=2(\ell(\mu)-(t+1-2))\ge 2(\ell(\mu)-(t+1-2))-1,\nonumber\\
&G(\ell(\mu)-(t+1-1),1)\le 2n-(2(t+1-2)+4)=2\ell(\mu)-(2(t+1-2)+4)=2(\ell(\mu)-t)-2\nonumber\\
&\ngeq 2(\ell(\mu)-t)-1
\end{align}
Note $\ell(\mu^{2n-2(t+1-2)})=\ell(\mu)-(t+1-2)-1)=\ell(\mu)-(t+1-1)=\ell(\mu)-t$. Furthermore, $t$ is the number of $>$ in \eqref{vip} before arriving to the flat sequence in \eqref{flat} which is also the number of the first $t$ non negative even numbers, $t=\#\{0,2,4,\dots,2(t+1-2)\}$.

\iffalse
Hence, $ G(\ell(\mu),1)=2n\ge 2n-1$, $ G(\ell(\mu)-1,1)=2n-2\ge 2(n-1)-1, \dots,$ $$G(\ell(\mu)-(t+1-2),1)=2\ell(\mu)-2(t+1-2)=2(\ell(\mu)-(t+1-2))\ge 2(\ell(\mu)-(t+1-2))-1,$$
 and $G(\ell(\mu)-(t+1-1),1)=$
$G(\ell(\mu)-t,1)\ngeq 2(\ell(\mu)-t)-1$ and
\fi
Thus $G_\mu(T)$ is not symplectic with bottom most symplectic violation in the cell $(\ell(\mu)-t,1)$.

\medskip
In the remaining  cases there exists at least one odd number among

$$T(n+1,1)\ge 2\times 1, T(n+2,1)\ge 2\times 2,\dots,
T(n+t-2),1)\ge 2\times (t-2),T(n+t-1),1)\ge 2\times (t-1).$$

This means, at least one of the even numbers $2,4,\dots, 2(t-1)$ is replaced by an odd number, that is, for some $1\le i\le t-1$, $2i$ is replaced by $2i+1$ and $2(i+1)$ is preserved in the list.
Indeed such  odd numbers do not violate the Sundaram condition. Each time we do this we glue two next flats subsequences in \eqref{vip} and reduce by one unity the number of $>$.  The  flat tail \eqref{flat} can be longer but the flat portion
$\ell(\mu^{2n-(2(t+1-2)+1)})=\ell(\mu^{2n-(2(t+1-2)+2)})=\ell(\mu^{2n-(2(t+1-2)+3)})=\ell(\mu^{2n-(2(t+1-2)+4)})$ is preserved. Therefore,  in   the first column of $G_\mu(T)$, from the bottom, the first symplectic violation   occurs among the bottom most $t+1$ cells $(\ell(\mu),1)$,  $(\ell(\mu)-1,1),\dots, (\ell(\mu)-t,1)$.

\iffalse
A situation where the bottom most failing is in cell $(\ell(\mu),1)$ of $G$, is the following

\emph{Case $n=\ell(\mu)+t\ge t+2\Rightarrow\ell(\mu)\ge 2$} in which case one has

$T(\ell(\mu)+1,1)=1$, $T(\ell(\mu)+2,1)=2$, $T(\ell(\mu)+3,1)=3$, $\dots,T(\ell(\mu)+t,1)=t$, $T((\ell(\mu)+t)+1,1)=t+1$, $\dots,T((\ell(\mu)+t)+t,1)=2t,$ $ T((\ell(\mu)+t)+t+1,1)=2t+1$, $T((\ell(\mu)+t)+t+2,1)=2(t+1)$.

This translate to a sequence where all $>$ in \eqref{vip} disappear and we are reduced to a longest final flat sequence in \eqref{flat}
\begin{align}&\ell(\mu^{2n})=\ell(\mu^{2n-1})=\dots=\ell(\mu^{2n-t+1})=\ell(\mu^{2n-t})
=\ell(\mu^{2n-t-1})=\ell(\mu^{2n-t-2})=\cdots \ge.\nonumber
\end{align}
\fi
\end{proof}

\end{document}
